\subsection{Exact relation-module bounds beyond the finite base}

Use the discrete group $U_0$ and its subgroup $L^6$
from B.2.15, and put $V_p=H_1(L,\mathbb F_p)$.
The action of $\mathcal T=B_0/L$ on this vector space
is specified by the Bass--Serre graph of the same
marked quotient $B_0\to\mathcal T$.  It has one edge
for each element of $\mathcal T$, with endpoints the
cosets of $\langle s\rangle$ and
$\langle r\rangle$.  In particular it has 144 edges,
120 vertices, and cycle rank 25.  Its exact integral
chain sequence is
\[
0\longrightarrow H_1(L,\mathbb Z)
 \longrightarrow\mathbb Z[\mathcal T]
 \longrightarrow
 \mathbb Z[\mathcal T/\langle s\rangle]
 \oplus\mathbb Z[\mathcal T/\langle r\rangle]
 \longrightarrow\mathbb Z\longrightarrow0.                    \tag{B.S38}
\]
All actions below come from left translation on this
graph; using inverse left translation gives the
equivalent right-module convention in the calculation.

\begin{result}{B.2.21}{Proved}{Generator and index bounds for the remaining obstruction}
Let $\overline M=M/S$ be the finite obstruction group
of B.2.19.  Its two-primary part is generated by at
most four elements, and its three-primary part is
cyclic.  Together with the annihilator in B.2.15,
this gives the explicit bounds
\[
 |\overline M|\leq 2^{56}3^6,
       \qquad [H:D_P(H)]\leq 2^{56}3^6.                       \tag{B.S39}
\]
These assertions bound the remaining classes; they
do not assert that any spectral-sequence term survives.
\end{result}

\begin{proof}
The group $E=U_0/L^6$ acts transitively on the six
free factors.  The stabilizer of the first is
$J_{\mathcal T}=J_0/L^6$, and its action on the first
factor's homology is its first-coordinate projection
$J_{\mathcal T}\to\mathcal T$.  Thus the actual
$E$-module of first homology is
\[
 H_1(L^6,\mathbb F_p)
       \cong\operatorname{Ind}_{J_{\mathcal T}}^E V_p.
                                                                  \tag{B.S40}
\]
Shapiro's chain isomorphism therefore identifies its
first $E$-homology with $H_1(J_{\mathcal T},V_p)$.
The marked graph (B.S38) and the exact power--conjugate
presentation of $J_{\mathcal T}$ give
\[
\begin{array}{c|c|c|c|c}
 p&\text{cocycle unknowns}&\text{relation rank}&
     \text{coboundary rank}&\dim H_1(J_{\mathcal T},V_p)\\ \hline
 2&450&424&25&1\\
 3&450&424&25&1
\end{array}                                                     \tag{B.S41}
\]
Here homology is calculated by dualizing to first
cohomology of the contragredient module; thus the
last entry is $450-424-25$.  For verification, the
same calculation on $\mathcal T$ itself has 150
unknowns, relation rank 124, and coboundary rank 25
at both primes.

We also retain the full second-homology coinvariants.
The Kunneth formula, now over $\mathbb F_p$, gives
\[
 H_2(L^6,\mathbb F_p)
                    =\bigoplus_{i<j} V_{p,i}\otimes V_{p,j}.
\]
The nine pair orbits with trivial common quotient
vanish after coinvariants: their action is the
independent product action, and
$(V_p)_{B_0}=0$ by $[B_0,L]=L$ in B.2.11.
For the two exceptional matchings, take coinvariants
by the independent kernels first and then use their
actual stabilizer actions from B.2.15.  The integral
graph action gives the following complete ranks:
\[
\begin{array}{c|c|c|c|c}
 p&\substack{\text{common}\\\text{quotient}}&
 \substack{\dim\text{ factor}\\\text{coinvariants}}&
 \substack{\text{tensor/stabilizer}\\\text{relation rank}}&
 \substack{\dim\text{ final}\\\text{coinvariants}}\\ \hline
 2&S_3&2&3&1\\
 2&A_4&3&7&2\\
 3&S_3&2&4&0\\
 3&A_4&3&9&0
\end{array}                                                     \tag{B.S42}
\]
The first stabilizer is the signed swap.  The second
is the signed swap followed by the actual $s$ action
on its second factor; the common-quotient generator
maps are exactly those in B.2.15.  Thus the dimensions
of $H_0(E,H_2(L^6,\mathbb F_p))$ are three and zero.

Equations (B.S41)--(B.S42) are reproduced by
\path{work/centralizer_code/export_universal_relation_module.g}
and
\path{work/centralizer_code/check_universal_relation_module.py}.
The exported data specify all 144 graph actions and
both exact finite presentations.  The Python calculation
checks the integral cycle boundaries, reconstructs every
translated cycle, verifies every presentation relator
on the module, and then computes the displayed ranks
over the stated finite fields.  The small tensor ranks
are independently checked by a second row-reduction
routine.  Its \texttt{.json} and \texttt{.out} files
retain the complete resulting tables.

Apply the homology spectral sequence to
$L^6\to U_0\to E$.  The three terms which can
contribute to total degree two have dimensions
\[
\begin{array}{c|c|c|c|c}
 p&\dim H_2(E,\mathbb F_p)&
 \dim H_1(E,H_1(L^6,\mathbb F_p))&
 \dim H_0(E,H_2(L^6,\mathbb F_p))&
 \text{sum}\\ \hline
 2&2&1&3&6\\
 3&2&1&0&3
\end{array}                                                     \tag{B.S43}
\]
The first column is (B.S36), since finite-dimensional
homology and cohomology with trivial field coefficients
are dual.  Each term of the limiting filtration is a
subquotient of its displayed term.  Therefore
$\dim H_2(U_0,\mathbb F_2)\leq6$ and
$\dim H_2(U_0,\mathbb F_3)\leq3$, regardless of
the differentials.

The integral abelianization of $U_0$ is $C_6$.
The universal-coefficient exact sequence consequently
subtracts one dimension at each prime from these
bounds for the reduction of $H_2(U_0,\mathbb Z)$.
The actual finite $M$ is a quotient of this integral
homology group by B.2.15.  Its elementary quotients
therefore have dimensions at most five and two.

Finally the two maps in (B.S34) lift to $S\subseteq M$.
Their images in $M$ have orders two and three and map
onto the respective factors of $M_E$.  Their product
is consequently a subgroup of $S$ mapping
isomorphically to $M_E\cong C_2\times C_3$.
Since $M$ is abelian, it gives an actual splitting
$M\cong M_0\times C_2\times C_3$.  It removes one
generator at each prime, so $M_0$ and its quotient
$\overline M$ have the asserted generator bounds.
Their exponent divides $2^{14}3^6$ by B.2.15, which
proves the first inequality in (B.S39).  The exact
orbit formula (B.S29) proves the second.
\end{proof}
